\pagestyle{plain} % No headers, just page numbers
% \pagenumbering{arabic} % Arabic numerals
% \setcounter{page}{1}


\chapter{\uppercase {Related Work}}

\section{Dialog Act Classification}
Several studies have explored different approaches to facilitate DA classification, with a focus on employing various neural network architectures and attention mechanisms for leveraging contextual information. \cite{ortega-vu-2017-neural} use CNNs to acquire utterance representations and further investigate various forms of attention mechanisms to infuse context awareness into the utterance representations. \cite{bothe-etal-2018-context} propose using a pre-trained character-level language model to generate sentence representations and a RNN to learn dialog act from the current utterance representation and the context of previous utterances. \cite{raheja-tetreault-2019-dialogue} propose a hierarchical RNN with context-aware self attention to model various levels of utterance and dialog act semantics. \cite{wang-etal-2020-integrating} propose Heterogeneous User History (HUH) graph convolution network, complemented by denoising mechanisms in order to effectively integrate user's historical answers grouped by DA labels to improve DA classification. \cite{he-etal-2021-speaker-turn} propose a speaker turn-aware approach where the turn embedding vector, learned based on the speaker label is combined with the utterance vector for further processing by an RNN to capture contextual information. \cite{zelasko-etal-2021-helps} explore the effects of broader context, capitalization and punctuation on DA classification and report strong result using XLNet and Longformer. 

Multimodality is less explored for DA classification.
\cite{he2018exploring} combine CNN for audio feature augmentation and RNN to model utterances from individual modalities and fuse them through concatenation. 
While \cite{lexico} utilize CNN to capture intra-utterance context for both audio and textual modalities, they only consider dialog-level context for modeling textual representation.
\cite{saha-etal-2020-towards} propose a Triplet Attention Subnetwork, incorporating self and cross attention to jointly model emotion and DA in dialogs.


\section{Dialog Breakdown Detection} The Dialog Breakdown Detection Challenge (DBDC) has been a pivotal platform for advancing research in this area~\cite{higashinaka-etal-2016-dialogue,HORI20191}. \cite{higashinaka-etal-2016-dialogue} defines the task description, datasets, and evaluation metrics for DBDC and provides insights into the design and methods used in these challenges. The methodologies employed range from traditional machine learning techniques to advanced neural network models. \cite{Hendriksen2021} explore different variants of LSTM for dialog breakdown detection. \cite{sugiyama2021dialogue} demonstrate a novel approach on dialog breakdown detection by integrating BERT's powerful language understanding capabilities with traditional dialog features like dialog acts. Another significant contribution is the exploration of semi-supervised learning methods to improve dialog breakdown detection, as discussed in \cite{Ng2020ImprovingDB}. Their research demonstrates the use of continued pre-training on the Reddit dataset and a manifold-based data augmentation method, showing a substantial improvement in detecting dialog breakdowns. The findings across these papers consistently indicate that the integration of advanced language models with contextual and conversational features significantly enhances the detection of dialog breakdowns. A few approaches use acoustic signals in previous related works. \cite{meena-etal-2015-automatic} used the output from automatic speech recognition system (ASR) systems as features to detect dialog breakdowns from spoken dialogs but they used only surface forms of STT texts or extracted text features from them rather than latent vectors of acoustic signals directly. 

\section{Deception Detection}

Early research on automated deception detection leveraged handcrafted linguistic, syntactic, and lexical features, including Linguistic Inquiry and Word Count (LIWC) indicators, part-of-speech distributions, and n-gram features, to capture linguistic, psychological and stylistic patterns indicative of deception. These features were utilized in statistical models such as logistic regression, decision trees, and support vector machines (SVM) to classify deceptive and truthful statements~\cite{opspam, rltd, levitan-etal-2018-linguistic, Ekman2001, LIWC3, LIWC4}. Audio-based deception detection has relied on Mel-frequency cepstral coefficients (MFCCs) and prosodic cues, such as pitch and speaking rate, to distinguish deceptive from truthful speech~\cite{Hirschberg2005DeceptiveSpeech, levitan-etal-2018-linguistic, Bai2019AutomaticLD, BoL, Chebbi2021DeceptionDU}. Additionally, research in nonverbal deception detection has focused on facial Action Units (AUs) extracted from video data, which capture microexpressions and facial muscle movements associated with deceptive behavior~\cite{Ekman2001, fau1, Bai2019AutomaticLD, LIWC3, mathur2021affectawaredeepbeliefnetwork}. These approaches, though effective in constrained settings, often struggle with generalization across datasets and speaker variations, necessitating the exploration of more robust deep learning techniques.

Recent advances in deep learning have led to an increasing adoption of CNNs and LSTMs for deception detection tasks across both textual and multimodal domains~\cite{DL5, DL2, DL4, DL3, DL1}. Transformer based models and attention mechanisms have also been applied in recent deception detection research, leveraging contextual embeddings to capture subtle deception cues~\cite{tr1, tr2, tr3}. ~\cite{dolos} presents a novel method called Parameter-Efficient Crossmodal Learning (PECL) that uses a temporal adapter to capture temporal attention and a fusion module to merge audio and visual cues for audio-visual deception detection. Building on these developments, emerging research is now harnessing the capabilities of LLMs—whose success across diverse cognitive tasks underscores their potential—to capture intricate linguistic nuances and further enhance deception detection. In their study, \cite{LLM1} employ variants of the FLAN-T5 model~\cite{Flant5} to detect deception across a range of textual contexts. \cite{LLM2} investigates the effectiveness of LLMs in deception detection using a Retrieval Augmented Generation (RAG) framework for few-shot learning in various textual domains. Our work advances this line of research by investigating the application of LLMs in real-world multimodal scenarios.